% !TeX root = ../papers/nc-p3.tex
\section{反対階数12の量子極限と正則性平面の極}
\label{sec:rank-twelve-quantum-limit}
\begin{theorem}\label{thm:rank-twelve-quantum-limit}
Pauli八係数の五つの斉次二次式の射影解のうち $a=b=0$ となるものは32点で、
全て同一の添字保持同型類の四窓を与える。その反対階数は12であり、
四周期の基底変更で次の量子多項式環の窓になる。
\[
 X_jX_k=q_{jk}X_kX_j\quad(j<k),\qquad
 (q_{01},q_{02},q_{03},q_{12},q_{13},q_{23})=(1,1,-i,-i,1,-1),
 \qquad i^2=-1 .
\]
従ってこの境界類は標準Koszul左右AS正則で、最小周期は1である。
\end{theorem}
\begin{proof}
$L^2+M^2=A^2+B^2=C^2+D^2=0$、$LM=-AB=-CD$ なので
$L=0$ なら全係数が零になる。$L=1$ とすると
$M=\pm i$, $A^4=C^4=1$, $B=-M/A$, $D=-M/C$。
\path{computations/verify_rank_twelve_quantum_limit.py} は192変換による
$(L,M,A,B,C,D)=(1,-i,i,1,i,1)$ の射影軌道がこの32点と一致することを検算する。
各変換は四位置の基底変更を与える。非零なPauli係数は十二個であり、
両側のPauli行列がそれぞれ基底をなすので反対切断の階数は12である。

代表の各文字の基底で対 $(0,1),(2,3)$ に
$\left(\begin{smallmatrix}1&1\\1&-1\end{smallmatrix}\right)$ を施し、
奇数位置では文字を $j\mapsto j\oplus3$ で付け替える。四次テンソルの
非零な24成分は共通倍率を除き $\pm1\pm i$ である。各巡回隣接像は
六つの二項式を与え、順序を定理と同じに取ると係数は
\[
\begin{array}{c|rrrrrr}
0&1&1&-i&-i&1&-1\\
1&1&-1&i&i&-1&-1\\
2&-1&-1&-i&-i&-1&1\\
3&-1&1&i&i&1&1
\end{array}
\]
となる。各位置 $p$ の対角基底倍率を $d_{p,j}$ とする。
隣接比 $d_{p+1,j}/d_{p,j}$ を順に
$(1,1,1,1)$、$(1,1,-1,-1)$、$(1,-1,-1,1)$、$(1,-1,1,-1)$ とすると、
全行が最初の行になり、各文字で四つの比の積は1である。
これらの計算は $2$ が単元の任意の剰余体でも成立する。
順序単項式は基底をなし、四乗 $X_j^4$ による商には各指数が4未満の単項式が残る。
各単項式 $X^e$ に沿う四周期の中心切断は
$\prod_j q_{jk}^{e_j}$ が $k$ によらないとき一次元、それ以外は零である
（逆順の係数は逆数とする）。共通値の四乗は1なので一周の整合も成立する。
次数4では四つの純四乗と $e=(1,1,1,1)$ の計五本を得る。
\end{proof}

\begin{theorem}\label{thm:rank-fifteen-pole-cone}
標数零の通常な階数15窓の正則性座標 $x,y$ に付値 $\nu$ を取り、
剰余標数は零または奇素数とする。$\nu(0)=+\infty$ として
\[
                 \nu(y)<\min\{0,2\nu(x)\}
\]
ならばこの窓は標準Koszul左右AS正則で、各始点で四中心四次式の
多項式環上自由階数256となる。
\end{theorem}
\begin{proof}
$D=x^3/y$ として $(U+3)^3=D(U-1)^2$ を解く。
$D$ が整なら全根は整で $U-1$ は単元である（$4$ が単元）。
$V^3=y(U-1)^4$ より $V$ は極を持つ。$D$ が極なら $d^2=1/D$ として
$U=1+dE$ と置く。方程式 $(4+dE)^3=E^2$ の剰余根 $E=\pm8$ は単根であり、
付値を延長して単元根を選べる。このとき
$3\nu(V)=\nu(y)-2\nu(D)=3\nu(y)-6\nu(x)<0$ である。
いずれも $U$ は整、$V$ は極である。$x\ne0$ では
$V=x(U-1)^2/(U+3)$ を取り、$x=0$ では $U=-3$ として上の三乗根を取る。

$K^2=U$, $H^2=V$ とし、$r+1/r=2(K+H)$、$w+1/w=2(K-H)$
で極を持つ方の根を取る。$s^2=r$, $t^2=w$ の符号を選んで
$\overline{t/s}=i$、$\overline{H/s^2}=1/2$ となる。
$n/s$ の整方程式の還元は $e^4=-1/4$ なので、単根 $e=(1+i)/2$ を指定できる。
八係数全体を $4/s$ 倍すると $a,b$ の剰余は零、
$(L,M,A,B,C,D)$ の剰余は
$(2e,1/e,i-1,i+1,i-1,i+1)$ である。
$1+i$ で割ると前定理の代表を得る。特に $K=2$ を固定して
$H$ を無限遠へ動かせば、通常な階数15からこの境界類への退化が実現する。

剰余の六次元の隣接像は持ち上がり、長さ4,5の経路次元は $35,56$、
四次中心の全次元は5で、いずれも一般点の下界に等しい。
従って剰余の四つの純四乗を中心切断に持ち上げられる。
剰余の四中心商の基底は上で与えたので、長さ13で消え、長さ0から12で
$(1+t+t^2+t^3)^4$ の係数を与える。一般点からの経路次元の下界と合わせて
全次数の自由性とKoszul性を得る。選んだ $(U,V)$ の表示と元の通常窓は
同じ $(x,y)$ を持つので、基底変更・巡回移動・反転により元の主張が従う。
\end{proof}

\begin{corollary}
$p=0$ で $y$ が代数的なら、未認証点の候補は、$y$ が全ての奇素数上で整であり、
かつ3の上の全ての素点で正の付値を持つ場合に限られる。
\end{corollary}
\begin{proof}
$x^3=27y$ なので奇素数上で $y$ が極なら前定理の厳密不等式を満たす。
3の上で $y$ が単元なら $\bar x=0$ で定理\ref{thm:rank-fifteen-residue-line}が適用できる。
\end{proof}
